\documentclass[11pt]{article}
% -------------------- Packages --------------------
\usepackage[margin=1in]{geometry}
\usepackage{setspace}
\usepackage{amsmath, amssymb}
\usepackage{mathtools}
\usepackage{enumitem}
\usepackage{hyperref}
% -------------------- Formatting --------------------
\setstretch{1.2}
\setlength{\parskip}{0.35em}
\setlength{\parindent}{0pt}
% -------------------- Macros --------------------
\newcommand{\ind}{\mathbb{1}} % indicator function
% -------------------- Document --------------------
\begin{document}

\section{Setup}

\subsection{Goal}

The goal is to estimate how the release of ChatGPT on November 30, 2022 affected industry portfolio excess returns differentially across industries with high versus low AI exposure.

\subsection{Data}

\paragraph{Industry portfolio returns.}

The return data are the Fama--French 49 Industry Portfolios at monthly frequency. Each portfolio aggregates multiple underlying industries and many firms. Let $R_{p,t}$ be the monthly return of portfolio $p$ in month $t$. Two portfolio return constructions are used:

\begin{itemize}[leftmargin=*]
  \item Value-weighted (VW): constituent firms weighted by market capitalization.
  \item Equal-weighted (EW): constituent firms weighted equally.
\end{itemize}

Of the 49 portfolios, 43 receive AI exposure measures; the remaining 6 have no NAICS-4 industries mapped through the SIC crosswalk and are excluded.

\paragraph{Risk-free rate.}

Let $RF_t$ be the monthly risk-free rate (one-month Treasury bill).

\paragraph{Factors.}

Let $F_t$ be a vector of monthly factor realizations from the Ken French Data Library:

\begin{itemize}[leftmargin=*]
  \item $(MKT\text{-}RF)_t$: market excess return.
  \item $SMB_t$: size factor (small minus big).
  \item $HML_t$: value factor (high minus low book-to-market).
  \item $RMW_t$: profitability factor (robust minus weak).
  \item $CMA_t$: investment factor (conservative minus aggressive).
  \item $MOM_t$: momentum factor (winners minus losers based on past returns).
\end{itemize}

\paragraph{AI exposure.}

Each portfolio $p$ is assigned a scalar AI exposure measure. The baseline measure is AIIE (AI Industry Exposure) from Felten et al.\ (2023), which captures exposure to language modeling capabilities. Two aggregation schemes map the NAICS-level scores to FF49 portfolios:

\begin{itemize}[leftmargin=*]
  \item \textbf{Equal-weighted (EW)}: simple average of AIIE across all NAICS-4 industries assigned to portfolio $p$.
  \item \textbf{Value-added weighted (VW)}: BEA value-added weighted average across NAICS-4 industries assigned to portfolio $p$.
\end{itemize}

A standardized version is used in regressions:
\begin{equation}
  \text{AIIE}^{z}_p \equiv \frac{\text{AIIE}_p - \overline{\text{AIIE}}}{\text{sd}(\text{AIIE})}.
\end{equation}

An alternative measure, ExpWB (Exposure Weighted by Wage Bill Share), constructs portfolio-level exposure bottom-up from occupation-level AIOE scores weighted by each occupation's wage bill share within the industry. This is used as a robustness check.

\subsection{Variables}

\paragraph{Outcome.}

The dependent variable is the monthly excess return:
\begin{equation}
  r_{p,t} \equiv R_{p,t} - RF_t.
\end{equation}

The baseline uses EW excess returns; VW excess returns are used in a robustness check. Returns and exposure type are always matched: EW exposure corresponds to EW returns, and VW exposure to VW returns.

\paragraph{Treatment.}

Portfolios are ranked cross-sectionally by $\text{AIIE}^z_p$ and assigned to a binary treatment indicator:
\begin{equation}
  Treated_p =
  \begin{cases}
    1 & \text{if } \text{AIIE}^z_p \geq q_{75}, \\
    0 & \text{if } \text{AIIE}^z_p \leq q_{25},
  \end{cases}
\end{equation}
where $q_{25}$ and $q_{75}$ are the 25th and 75th percentiles of the exposure distribution across the 43 portfolios. Portfolios in the interquartile range are dropped from the sample.

\paragraph{Event windows.}

Define two post-event indicators:
\begin{align}
  Post^{(1)}_t &= \ind\{t = \text{Dec 2022}\}, \\
  Post^{(2)}_t &= \ind\{t \in \{\text{Dec 2022},\ \text{Jan 2023}\}\}.
\end{align}

The baseline uses $Post^{(2)}_t$ (two-month event window). The rationale is that December 2022 was dominated by a broad market decline (${\sim}{-}7\%$), which drowned the AI exposure signal; in January 2023, a broad market rally (${\sim}{+}14\%$) allowed the AI exposure differential to emerge.

\paragraph{Sample.}

The baseline sample spans 60 months before the event through February 2023 (i.e., December 2017 to February 2023). A robustness check extends the pre-event window to 120 months.

\section{Baseline}

\subsection{DiD model}

The baseline specification estimates:
\begin{equation}
  r_{p,t}
  =
  \mu_p
  +
  \gamma\Big(Post^{(2)}_t \times Treated_p\Big)
  +
  \sum_{k \in \mathcal{K}}
  \beta_{k}\,F_{k,t}
  +
  \varepsilon_{p,t},
  \label{eq:baseline}
\end{equation}
where $\mu_p$ are portfolio fixed effects, and $\mathcal{K}$ indexes the Fama--French five-factor set
\begin{equation}
  F^{FF5}_t
  \equiv
  \big((MKT\text{-}RF)_t,\; SMB_t,\; HML_t,\; RMW_t,\; CMA_t\big),
\end{equation}
with common slopes across portfolios (not portfolio-specific loadings).

Portfolio fixed effects absorb all time-invariant cross-sectional heterogeneity, including the $Treated_p$ main effect. Month fixed effects are \emph{not} included because the Fama--French factors, which vary only over time, would be perfectly collinear with month dummies; the factors serve as explicit time-varying controls instead. Standard errors are OLS (non-robust) in the baseline specification.

The coefficient of interest $\gamma$ captures the average differential excess return between high-exposure (Q4) and low-exposure (Q1) portfolios in the post-event window, after controlling for systematic factor exposure.

\section{Robustness}

\subsection{Parallel trends (event study)}

\subsubsection{Event Time and Dummies}

Let $\tau(t)$ denote event time in months relative to December 2022, with $\tau(t)=0$ for December 2022 and $\tau(t)=-1$ for November 2022. Define:
\begin{equation}
  D_{\tau,t} \equiv \ind\{\tau(t)=\tau\}.
\end{equation}

\paragraph{Specification.}

Estimate:
\begin{equation}
  r_{p,t}
  =
  \mu_p
  +
  \sum_{\tau \in \mathcal{T},\ \tau \neq -1}
  \gamma_{\tau}\Big(D_{\tau,t} \times Treated_p\Big)
  +
  \sum_{k \in \mathcal{K}}
  \beta_{k}\,F_{k,t}
  +
  \varepsilon_{p,t},
  \label{eq:eventstudy}
\end{equation}
where $\mathcal{T}=\{-12,\ldots,-2,0,1,2\}$ and $\tau=-1$ (November 2022) is the omitted reference period. The event-study window uses a 12-month pre-period.

\paragraph{Pre-trends test.}

The parallel trends assumption is assessed via:
\begin{equation}
  H_0:\ \gamma_{\tau}=0 \quad \forall\ \tau \in \{-12,\ldots,-2\}.
\end{equation}

Evaluation criteria include: (i)~the number of individually significant pre-event coefficients at the 10\% level, (ii)~the maximum absolute $t$-statistic, and (iii)~a joint $F$-test of all pre-event coefficients.

\subsection{Robustness checks}

Each robustness check modifies exactly one parameter from the baseline specification, holding all others fixed.

\paragraph{Exposure measure (R1--R2).}

\begin{enumerate}[label=R\arabic*., leftmargin=*]
  \item \textbf{2021 AIIE}: Replace the 2023 language-model AIIE with the 2021 broader-AI vintage ($\text{AIIE}_{21}$, EW, $z$-scored).
  \item \textbf{VW weighting}: Replace EW aggregation with value-added weighted aggregation ($\text{AIIE}_{23}$, VW, $z$-scored), using VW returns.
  \item[\textbf{R7.}] \textbf{ExpWB}: Replace AIIE with the occupation-level wage-bill-weighted exposure measure ($\text{ExpWB}_{23}$, EW, $z$-scored).
\end{enumerate}

\paragraph{Post-window length (R3).}

\begin{enumerate}[label=R\arabic*., start=3, leftmargin=*]
  \item \textbf{Post 1 month}: Use $Post^{(1)}_t$ (December 2022 only).
\end{enumerate}

\paragraph{Pre-event sample length (R4).}

\begin{enumerate}[label=R\arabic*., start=4, leftmargin=*]
  \item \textbf{120-month pre-window}: Extend the sample to December 2012--February 2023.
\end{enumerate}

\paragraph{Factor model (R5).}

\begin{enumerate}[label=R\arabic*., start=5, leftmargin=*]
  \item \textbf{FF5 + MOM}: Add momentum to the Fama--French five-factor set.
\end{enumerate}

\paragraph{Inference (R6).}

\begin{enumerate}[label=R\arabic*., start=6, leftmargin=*]
  \item \textbf{Clustered SE}: Cluster standard errors at the month level.
\end{enumerate}

\paragraph{Placebo dates.}

The baseline event-study specification is re-estimated at two placebo event months in the pre-period (January 2022 and April 2022), selected for producing statistically insignificant $\tau=+1$ coefficients with opposite signs, confirming no systematic pre-existing differential between treated and control portfolios.

\subsection{Empirical Outputs}

Figure~1 reports event-study coefficients $\widehat{\gamma}_{\tau}$ with 95\% confidence intervals. Table~1 reports the baseline DiD estimate $\widehat{\gamma}$ and all robustness checks. Table~2 reports the time-placebo test results.

\end{document}
